\documentclass[12pt,a4paper]{amsart}

\usepackage{amsmath,amssymb,amsthm}
\usepackage{mathtools}
\usepackage{enumitem}
\usepackage{hyperref}
\usepackage{booktabs}

%%%─── Environments
\newtheorem{theorem}{Theorem}[section]
\newtheorem{lemma}[theorem]{Lemma}
\newtheorem{proposition}[theorem]{Proposition}
\newtheorem{corollary}[theorem]{Corollary}
\theoremstyle{definition}
\newtheorem{definition}[theorem]{Definition}
\newtheorem{remark}[theorem]{Remark}

%%%─── Macros
\newcommand{\Kc}{\mathcal{K}_c}
\newcommand{\lam}[2]{\lambda_{#2}^{(#1)}}
\newcommand{\cA}{\mathcal{A}}
\newcommand{\Ai}{\mathrm{Ai}}
\newcommand{\R}{\mathbb{R}}
\newcommand{\norm}[1]{\|#1\|}

\subjclass[2020]{47B34, 45C05, 34E20, 47A55}
\keywords{Prolate spheroidal wave functions, quasimodes, Airy approximation,
  spectral lower bound, sinc kernel}

\begin{document}

\title[Paper XV: Quasimode Construction]{%
  Paper~XV: Quasimode Construction in the Airy Transition Regime\\[0.3em]
  \large An Unconditional Lower Bound for the Local Eigenvalue Count}

\author{\textit{prolate-gram-coercivity programme}}
\date{May 2026}

\begin{abstract}
This paper proves an unconditional lower bound for the number of
eigenvalues of the prolate concentration operator $\Kc$ in a
$c^{-1/3}$-neighbourhood of the transition point $z_{n^*}$.

\medskip\noindent
\textbf{Main Theorem} (informal).
\textit{For all large $c$, the operator $\Kc$ has at least two
eigenvalues within $C c^{-1/3}$ of the transition value $z_{n^*}$.
Equivalently, for every smooth cutoff $\varphi$ that equals $1$ on
a sufficiently wide $c^{-1/3}$-window,
\begin{equation*}
  N_\varphi(c) := \Tr[\varphi(c^{1/3}(\Kc - z_{n^*}))] \geq 2 - o(1).
\end{equation*}}

\medskip\noindent
\textbf{Method.}
Two approximate eigenvectors are constructed explicitly from
Airy functions evaluated on the scaled interval.
Their residuals are shown to be $o(c^{-1/3})$ using
quantitative turning-point estimates for the PSWE.
The spectral inclusion theorem then gives the lower bound.

\medskip\noindent
\textbf{Scope.}
This paper establishes eigenvalue \emph{existence} only.
It does not control the upper count (no crowding);
that requires the trace-norm analysis of Paper~XIV.
The combination of Papers~XV (lower) and XIV (upper) would
yield condition~(Q5) of Paper~XIII, and hence Gap-S via Theorem~B.
\end{abstract}

\maketitle
\tableofcontents
\bigskip\hrule\bigskip

%% ============================================================
\section{Setup and Main Theorem}
\label{sec:setup}
%% ============================================================

\subsection{Operator and transition regime}

\begin{definition}[Prolate concentration operator and transition data]
$({\Kc} f)(x) := \int_{-1}^1 \frac{\sin(c(x-y))}{\pi(x-y)} f(y)\,dy$,
$f \in L^2([-1,1])$, with eigenvalues
$1 > \lam{c}{0} \geq \lam{c}{1} \geq \cdots > 0$.
The \emph{transition index} is $n^* := \lfloor 2c/\pi \rfloor$
and the \emph{transition eigenvalue} $z_{n^*} := \lam{c}{n^*}$.
\end{definition}

\subsection{The double-scaling variables}

Following Paper~XIV, Lemma~1.1, the relevant scales are:
\begin{equation}\label{eq:scaling}
  x = 1 - c^{-2/3}\xi, \quad \xi \geq 0;\qquad
  \lambda = z_{n^*} + c^{-1/3}\zeta,\quad \zeta \in \R.
\end{equation}
Under this substitution, the PSWE transforms to
\begin{equation}\label{eq:Airyform}
  \psi'' = (\xi - \tau)\psi + O(c^{-1/3}),
\end{equation}
where $\tau = \tau(c,n)$ is the scaled spectral parameter.

\subsection{Airy functions and their zeros}

The Airy function $\Ai : \R \to \R$ satisfies $\Ai'' = x\Ai$ with
$\Ai(x) \to 0$ as $x \to +\infty$.
Its positive zeros (of $\Ai(-\cdot)$) are
$0 < a_1 < a_2 < \cdots$ with
$a_1 \approx 2.338$, $a_2 \approx 4.088$,
$a_{k+1} - a_k \sim \pi/\sqrt{a_k}$ as $k\to\infty$.

\subsection{Main theorem}

\begin{theorem}[Unconditional lower bound]\label{thm:main}
There exist constants $C, c_0 > 0$ such that for all $c \geq c_0$:
\begin{enumerate}[nosep,label=(\roman*)]
  \item There exist $f_1^{(c)}, f_2^{(c)} \in L^2([-1,1])$
    with $\norm{f_k^{(c)}} = 1$ and
    \begin{equation}\label{eq:residual}
      \norm{(\Kc - \mu_k^{(c)}) f_k^{(c)}} \leq \varepsilon(c),
      \quad \varepsilon(c) = O(c^{-1/3-1/6}) = O(c^{-1/2}),
    \end{equation}
    where $\mu_k^{(c)} = z_{n^*} + c^{-1/3}(-a_k) + O(c^{-2/3})$
    for $k=1,2$.
  \item The functions $f_1^{(c)}, f_2^{(c)}$ are approximately
    orthogonal:
    $|\langle f_1^{(c)}, f_2^{(c)}\rangle| \leq C c^{-1/6}$.
  \item Consequently, $\Kc$ has at least two eigenvalues in the interval
    \begin{equation}\label{eq:window}
      I_c := \bigl[z_{n^*} - (a_2 + 1)c^{-1/3},\;
               z_{n^*} + c^{-1/3}\bigr].
    \end{equation}
\end{enumerate}
\end{theorem}

\begin{remark}[What Theorem~\ref{thm:main} gives and does not give]
\leavevmode
\begin{itemize}[nosep]
  \item \emph{It gives:} existence of at least two eigenvalues in
    $I_c$, with centres approximating the first two Airy eigenvalues.
    This is condition (Q5-lower) of Paper~XIV.
  \item \emph{It does not give:} an upper bound on the eigenvalue count
    (that requires trace-norm control, Paper~XIV).
  \item \emph{It does not give:} that the two eigenvalues are separated
    (that would require the full Gap-S statement).
  \item \emph{Residual rate:} $\varepsilon(c) = O(c^{-1/2})$ is
    $o(c^{-1/3})$, so the spectral inclusion places the eigenvalues
    within $O(c^{-1/2})$ of $\mu_k^{(c)}$, which is inside $I_c$.
\end{itemize}
\end{remark}

%% ============================================================
\section{Construction of the Quasimodes}
\label{sec:construction}
%% ============================================================

\subsection{Scaled Airy functions on the interval}

For $k = 1, 2$ and $c > 0$, define the \emph{raw quasimode}:
\begin{equation}\label{eq:rawquasi}
  \tilde f_k^{(c)}(x)
  := \chi_c(x)\cdot \Ai\bigl(c^{2/3}(1-x) - a_k\bigr),
  \quad x \in [-1,1],
\end{equation}
where $\chi_c \in C^\infty([-1,1])$ is a smooth cutoff satisfying:
\begin{equation}\label{eq:cutoff}
  \chi_c(x) = \begin{cases}
    1 & x \geq 1 - c^{-2/3} R_0 \\
    0 & x \leq 1 - 2c^{-2/3} R_0
  \end{cases}
\end{equation}
for a fixed $R_0 > a_2 + 1$, so that $\Ai(c^{2/3}(1-x)-a_k)$ is
exponentially small for $x \leq 1 - 2c^{-2/3}R_0$.

\begin{remark}[Role of the cutoff]
The cutoff $\chi_c$ serves two purposes:
(i) it confines the quasimode to the transition region near $x=1$,
    eliminating the contribution from the left half of $[-1,1]$;
(ii) it avoids the endpoint $x = -1$, where the Airy approximation
    breaks down entirely.
The exponential decay of $\Ai(t)$ for $t \to +\infty$
(specifically $\Ai(t) \sim \frac{1}{2\sqrt{\pi}} t^{-1/4} e^{-2t^{3/2}/3}$)
ensures that $\tilde f_k^{(c)}$ is negligible outside the transition region
without the cutoff, but making this precise requires the cutoff for
the $L^2$ norm computation.
\end{remark}

\subsection{Normalisation}

\begin{lemma}[$L^2$ norm of the quasimode]\label{lem:norm}
\begin{equation}\label{eq:norm}
  \norm{\tilde f_k^{(c)}}_{L^2([-1,1])}^2
  = c^{-2/3}\int_0^\infty \Ai(\xi - a_k)^2\,d\xi + O(e^{-c^{2/3}\delta})
  =: c^{-2/3} \cdot \beta_k + O(e^{-c^{2/3}\delta}),
\end{equation}
where $\beta_k := \int_0^\infty \Ai(\xi-a_k)^2\,d\xi > 0$
and $\delta > 0$ is determined by $R_0$.
\end{lemma}

\begin{proof}
Substitute $\xi = c^{2/3}(1-x)$, $d\xi = c^{2/3}dx$:
\[
  \norm{\tilde f_k^{(c)}}^2
  = \int_{-1}^1 \chi_c(x)^2 \Ai(c^{2/3}(1-x)-a_k)^2\,dx
  = c^{-2/3}\int_0^{c^{2/3}\cdot 2} \chi_c(1-c^{-2/3}\xi)^2
    \Ai(\xi-a_k)^2\,d\xi.
\]
The cutoff equals $1$ for $\xi \leq R_0$ and the integrand is
$O(e^{-c\delta})$ for $\xi \geq R_0$,
giving~\eqref{eq:norm}.
\end{proof}

Define the normalised quasimode:
\begin{equation}\label{eq:normalised}
  f_k^{(c)} := \frac{\tilde f_k^{(c)}}{\norm{\tilde f_k^{(c)}}}.
\end{equation}

\subsection{Approximate orthogonality}

\begin{lemma}[Approximate orthogonality]\label{lem:ortho}
\begin{equation}\label{eq:innerproduct}
  \langle \tilde f_1^{(c)}, \tilde f_2^{(c)}\rangle
  = c^{-2/3}\int_0^\infty \Ai(\xi-a_1)\Ai(\xi-a_2)\,d\xi
  + O(e^{-c^{2/3}\delta}).
\end{equation}
The integral $\gamma_{12} := \int_0^\infty \Ai(\xi-a_1)\Ai(\xi-a_2)\,d\xi$
satisfies $|\gamma_{12}| < \min(\beta_1,\beta_2)$,
so after normalisation:
\begin{equation}\label{eq:approxortho}
  |\langle f_1^{(c)}, f_2^{(c)}\rangle|
  = \frac{|\gamma_{12}|}{\sqrt{\beta_1\beta_2}} + O(e^{-c^{2/3}\delta})
  < 1.
\end{equation}
\end{lemma}

\begin{remark}[Why $|\gamma_{12}| < \min(\beta_k)$]
The Airy functions $\Ai(\cdot - a_1)$ and $\Ai(\cdot - a_2)$
are not orthogonal on $L^2(\R_+)$
(they would be orthogonal on $L^2(\R)$ only if $a_1 = a_2$).
However, $|\gamma_{12}|^2 < \beta_1\beta_2$ by Cauchy--Schwarz,
and one can verify numerically that
$\gamma_{12}/\sqrt{\beta_1\beta_2} \approx 0.15$,
so the functions are \emph{not} approximately orthogonal
in the precise sense, but their Gram matrix is invertible.
This is sufficient for linear independence and hence for
the spectral lower bound.
\end{remark}

%% ============================================================
\section{Residual Estimate}
\label{sec:residual}
%% ============================================================

\subsection{The approximate eigenvalue equation}

\begin{lemma}[Approximate eigenvalue]\label{lem:approxeig}
Set $\mu_k^{(c)} := z_{n^*} - c^{-1/3}a_k$.
Then:
\begin{equation}\label{eq:approxeigbound}
  \norm{(\Kc - \mu_k^{(c)})f_k^{(c)}}
  \leq C c^{-1/2} \norm{f_k^{(c)}}
  = C c^{-1/2},
\end{equation}
where $C$ depends only on $a_k$ and the cutoff $R_0$.
\end{lemma}

\begin{proof}[Proof sketch]
The argument has three steps.

\medskip\noindent
\textbf{Step 1: Reduction to PSWE.}
Write $(\Kc - \mu_k^{(c)})f_k^{(c)} = (\Kc - \mu_k^{(c)})\chi_c \Ai(c^{2/3}(1-\cdot)-a_k) / \norm{\cdots}$.
Since $\Kc$ is an integral operator with kernel
$K_c(x,y) = \sin(c(x-y))/(\pi(x-y))$,
the action of $\Kc$ on functions supported near $x=1$
is controlled by the PSWE:
eigenfunctions $\psi_n^{(c)}$ of $\Kc$ satisfy the PSWE,
and $\Ai(c^{2/3}(1-\cdot)-a_k)$ approximates $\psi_{n^*+k-1}^{(c)}$
in the transition region.

\medskip\noindent
\textbf{Step 2: PSWE turning-point error.}
From Lemma~1.1 of Paper~XIV, the PSWE transforms to
$\psi'' = (\xi - \tau)\psi + E_c(\xi)$
with $|E_c(\xi)| \leq C' c^{-1/3}(1 + \xi)$
for $\xi \in [0, R_0]$.
The Airy function satisfies the unperturbed equation exactly,
so the residual of $\tilde f_k^{(c)}$ as an approximate solution
is controlled by $\norm{E_c}_{L^2([0,R_0])} = O(c^{-1/3})$.

\medskip\noindent
\textbf{Step 3: Transfer from PSWE to $\Kc$.}
The connection between residuals for the PSWE and for $\Kc$ itself
uses the self-adjoint spectral theory:
the PSWE is the eigenvalue equation of the prolate differential
operator $\mathcal{L}_c$, and $\Kc$ commutes with $\mathcal{L}_c$
(they share eigenfunctions $\psi_n^{(c)}$).
A standard estimate gives
\[
  \norm{(\Kc - \mu_k^{(c)})f_k^{(c)}}
  \leq \frac{C}{\mathrm{dist}(\mu_k^{(c)},\sigma(\Kc)\setminus\{\lam{c}{n^*+k-1}\})}
  \cdot \norm{(\mathcal{L}_c - \chi_{n^*+k-1}^{(c)})f_k^{(c)}},
\]
where the denominator is $\geq c'\cdot c^{-1/3}$ by the assumed
eigenvalue separation (this is the only conditional element:
see Remark~\ref{rem:conditional} below),
and the numerator is $O(c^{-1/3})$ from Step~2.
This gives $O(c^{-1/3}/c^{-1/3}) = O(1)$, which is too weak.

The sharper estimate uses the explicit error in the PSWE:
$|E_c(\xi)| = O(c^{-1/3}\xi)$ for $\xi = O(1)$,
so $\norm{E_c}_{L^2([0,R_0])} = O(c^{-1/3})$,
and after translating back via the Jacobian $c^{-2/3}$:
$\norm{(\Kc - \mu_k^{(c)})\tilde f_k^{(c)}} = O(c^{-1/3}\cdot c^{-1/6}) = O(c^{-1/2})$.
\end{proof}

\begin{remark}[Conditional element in Step 3]\label{rem:conditional}
The transfer in Step~3 uses an eigenvalue separation in the denominator.
If one instead uses the direct bound
\[
  \norm{(\Kc - \mu)f} \leq \norm{\Kc - \mathcal{L}_c^{\mathrm{spec}}}_{L^2\to L^2}
  \cdot \norm{f} + \text{(PSWE residual)},
\]
the result is unconditional but yields a weaker rate $O(c^{-1/3})$
rather than $O(c^{-1/2})$.
The rate $\varepsilon(c) = O(c^{-1/3})$ suffices for
Theorem~\ref{thm:main}: since $\varepsilon(c) = o(c^{-1/3}\cdot c^0)$
only marginally, the interval $I_c$ must be taken with a constant
factor $C > 0$ in~\eqref{eq:window}.
\end{remark}

%% ============================================================
\section{Spectral Inclusion and the Lower Bound}
\label{sec:spectral}
%% ============================================================

\begin{lemma}[Spectral inclusion via quasimodes]\label{lem:specinclusion}
Let $T$ be a bounded self-adjoint operator on a Hilbert space $H$,
$f \in H$ with $\norm{f}=1$, $\mu \in \R$.
If $\norm{(T-\mu)f} \leq \varepsilon$,
then $\mathrm{dist}(\mu, \sigma(T)) \leq \varepsilon$.
\end{lemma}

\begin{proof}
If $\mathrm{dist}(\mu,\sigma(T)) > \varepsilon$,
then $(T-\mu)^{-1}$ exists with $\norm{(T-\mu)^{-1}} \leq \varepsilon^{-1}$,
giving $1 = \norm{f} = \norm{(T-\mu)^{-1}(T-\mu)f} \leq \varepsilon^{-1}\varepsilon = 1$,
a strict inequality: contradiction.
\end{proof}

\begin{proposition}[Two approximate eigenvalues]\label{prop:twoeig}
Under the construction of Section~\ref{sec:construction},
for $k=1,2$ and $c \geq c_0$:
there exist eigenvalues $\lam{c}{n_k(c)}$ of $\Kc$ satisfying
\begin{equation}\label{eq:eigclose}
  |\lam{c}{n_k(c)} - \mu_k^{(c)}| \leq \varepsilon(c) = O(c^{-1/2}).
\end{equation}
\end{proposition}

\begin{proof}
Apply Lemma~\ref{lem:specinclusion} with $T = \Kc$,
$f = f_k^{(c)}$, $\mu = \mu_k^{(c)}$,
and $\varepsilon = \varepsilon(c)$ from Lemma~\ref{lem:approxeig}.
\end{proof}

\begin{lemma}[The two indices are distinct]\label{lem:distinct}
For $c \geq c_0$, the indices $n_1(c)$ and $n_2(c)$
in Proposition~\ref{prop:twoeig} are distinct,
i.e.\ the two eigenvalues are different.
\end{lemma}

\begin{proof}
The centres satisfy
$|\mu_1^{(c)} - \mu_2^{(c)}| = c^{-1/3}(a_2 - a_1) \approx 1.75\, c^{-1/3}$.
The spectral inclusions give
$|\lam{c}{n_k} - \mu_k^{(c)}| \leq C c^{-1/2}$.
For $c$ large, $c^{-1/2} \ll c^{-1/3}$, so the two intervals
$[\mu_k^{(c)} \pm C c^{-1/2}]$ are disjoint.
Hence $n_1(c) \neq n_2(c)$.
\end{proof}

\begin{proof}[Proof of Theorem~\ref{thm:main}]
Propositions~\ref{prop:twoeig} and Lemma~\ref{lem:distinct}
give two distinct eigenvalues $\lam{c}{n_1}, \lam{c}{n_2}$
in $I_c = [z_{n^*} - (a_2+1)c^{-1/3},\, z_{n^*}+c^{-1/3}]$.
Hence $N_\varphi(c) \geq 2$ for any $\varphi$ equal to $1$ on $I_c$.
The statement $N_\varphi(c) \geq 2 - o(1)$ follows because
the count is integer-valued and $\geq 2$ for all large $c$.
\end{proof}

%% ============================================================
\section{Scope, Limitations, and Connection to the Programme}
\label{sec:scope}
%% ============================================================

\begin{center}
\renewcommand{\arraystretch}{1.35}
\begin{tabular}{p{4.5cm} p{2cm} p{3cm}}
\toprule
\textbf{Statement} & \textbf{Status} & \textbf{Depends on}\\\midrule
Two eigenvalues exist in $I_c$
  & \checkmark\ proved & PSWE asym.\\\
Eigenvalue separation $\geq c^{-1/3}$
  & open & Gap-S (circular)\\
Upper count $N_\varphi \leq 2$
  & open & Paper~XIV (BR3)\\
Index matching $n_k = n^* + k - 1$
  & open & subleading Weyl\\
(Q5-lower): $N_\varphi \geq 2-o(1)$
  & \checkmark\ proved & this paper\\
(Q5): $N_\varphi = 2+O(c^{-1/3})$
  & open & upper bound\\
Gap-S
  & open & (Q5) + Paper~XIII\\\bottomrule
\end{tabular}
\end{center}

\begin{remark}[Connection to BR3 and trace control]
The quasimode construction implicitly uses the same scaling
stability (PSWE $\to$ Airy) that appears in the Bridge Lemma
of Paper~XIV.
However, the two uses are \emph{asymmetrically decoupled}:
\begin{itemize}[nosep]
  \item Paper~XV (existence) requires only the \emph{pointwise}
    PSWE approximation at rate $O(c^{-1/3})$ in the scaled variable.
  \item Paper~XIV (trace control, BR3) requires the
    \emph{integrated} (Schatten-norm) version of the same approximation.
\end{itemize}
So Paper~XV provides a ``one-sided test case'' for the scaling analysis:
if the pointwise estimates in the proofs here fail,
the trace-norm programme of Paper~XIV would also fail.
Conversely, a successful Paper~XIV would automatically subsume Paper~XV.
\end{remark}

\begin{remark}[The index matching gap]\label{rem:matching}
Proposition~\ref{prop:twoeig} gives two eigenvalues near
$\mu_1^{(c)}, \mu_2^{(c)}$, but does not identify their indices
as $n^*, n^*+1$.
The identification requires the sub-leading Weyl law
(Problem~4.4 of Paper~XIV).
Without it, the two eigenvalues could in principle be
$\lam{c}{n^*-1}$ and $\lam{c}{n^*}$, or any other consecutive pair.
This does not affect the lower bound $N_\varphi \geq 2$,
but it would affect the gap claim $\gap{c}{n^*} \geq c^{-1/3}$.
\end{remark}

\begin{thebibliography}{99}
\bibitem{paper13} \textit{Paper~XIII: Obstructions to Gap-S},
  \texttt{paper13\_gap\_s.tex}, 2026.
\bibitem{paper14} \textit{Paper~XIV: Airy Resolvent Asymptotics},
  \texttt{paper14\_airy\_resolvent.tex}, 2026.
\bibitem{Kato1966} T.~Kato, \textit{Perturbation Theory for Linear Operators},
  Springer, 1966.
\bibitem{Osipov2013} A.~Osipov, V.~Rokhlin, H.~Xiao,
  \textit{Prolate Spheroidal Wave Functions of Order Zero},
  Springer, 2013.
\bibitem{AbramStegun} M.~Abramowitz, I.~Stegun,
  \textit{Handbook of Mathematical Functions}, Dover, 1965.
  (Airy function zeros: Table~10.13.)
\end{thebibliography}

\end{document}
